\chapter{Crypto Data and Infrastructure}\label{ch:m3:crypto-data-and-infrastructure}

A purchased history of a crypto order book shows, on a quiet afternoon, the best bid above the best offer
for two seconds. Nothing happened in the market. One update of the venue's websocket feed was lost on the
way to the recorder, the recorder applied the next ones anyway, and every system downstream, a backtest,
a research notebook, a model trained on the data, treated the crossed book as real. Crypto data looks easy
to get: the venues publish their books for free, the chains are public. It is hard to get right. This
chapter covers the two sources of a crypto firm's data, the chains and the venues: running nodes and
reading what contracts emit, rebuilding venues' books from snapshots and deltas and checking them,
understanding what each timestamp measures, and knowing where the matching engines physically are.

\section{Running nodes}

\begin{definition}[Full node, archive node, RPC endpoint]\label{def:m3:crypto-data-and-infrastructure:node}
A \emph{full node}\index{full node} is a program that downloads and verifies every block of a chain and keeps its current
state, and the recent history needed to serve it. An \emph{archive node}\index{archive node} is a full node that also keeps
every historical state, so that it can answer questions about any past block. An \emph{RPC endpoint}\index{RPC endpoint} is the
interface through which a program queries a node and submits transactions to it.
\end{definition}

A trading firm reads the chain through nodes: its own, or a provider's. Its own nodes cost hardware and operations but give it
the chain's data as soon as its node sees it, no \omterm{def:m3:centralised-exchanges:ratelimit}{rate limits} and no dependence on a third party; a provider's endpoints are
convenient and are shared with everyone else, throttled, and one more party between the firm and the chain. For latency-sensitive
work (\cref{ch:m3:maximal-extractable-value}) the firm runs its own nodes, close to the other nodes it needs to hear from first.

\begin{dated}{2026-09}{Ethereum nodes}\label{dat:m3:crypto-data-and-infrastructure:nodes}
Ethereum's developer documentation states that \omterm{def:m3:crypto-data-and-infrastructure:node}{full nodes} verify every block but keep only relatively recent state (typically the
last 128 blocks), regenerating older data on demand; that \omterm{def:m3:crypto-data-and-infrastructure:node}{archive nodes} never delete downloaded data and represent terabytes of
storage; and that one client, Erigon, can perform a full archive sync using around 2~TB of disk in under three days. It lists five
execution clients (Geth, Nethermind, Besu, Erigon, Reth) written in four languages.
\end{dated}

\section{On-chain data: logs and indexers}

\begin{definition}[Event log, indexer]\label{def:m3:crypto-data-and-infrastructure:log}
An \emph{event log}\index{event log} is a record a \omterm{def:m3:blockchains-for-traders:contract}{smart contract} emits during a transaction (a swap, a transfer, a liquidation),
stored in the transaction's receipt with indexed fields that can be searched. An \emph{indexer}\index{indexer} is a system that reads
every block from a node, decodes the logs and state changes of the contracts it cares about, and stores them in a database that can
be queried by time, address or event type.
\end{definition}

The chain is a database with an unhelpful interface: it answers ``what is the state now'' and ``what happened in block $n$'', but not
``every swap in this pool last month''. The firm's \omterm{def:m3:crypto-data-and-infrastructure:log}{indexer} turns the one into the other. Its traps are those of any data pipeline with
a twist: blocks can be reorganised away (\cref{ch:m3:blockchains-for-traders}), so recent data must carry its block hash and be revised
when a reorganisation happens; contract upgrades change what logs mean; and a \omterm{def:m3:blockchains-for-traders:key}{token}'s decimals, not its log, say what a raw amount is
worth.

\begin{omfigure}
\centering
\begin{tikzpicture}[font=\small,
  s/.style={draw,thick,rounded corners=2pt,align=center,minimum height=0.9cm,text width=2.3cm},
  arr/.style={-{Stealth[length=2.2mm]},thick}]
\node[s,draw=omDef,fill=omDef!8] (ch) at (0,1.1) {chain\\(own nodes)};
\node[s,draw=omThm,fill=omThm!8] (ix) at (3.2,1.1) {indexer:\\logs, state};
\node[s,draw=omDef,fill=omDef!8] (vx) at (0,-1.1) {venues'\\websockets};
\node[s,draw=omThm,fill=omThm!8] (bb) at (3.2,-1.1) {book builder:\\ids, checksums};
\node[s,draw=omProp,fill=omProp!8] (st) at (6.6,0) {normalised\\store, two\\timestamps};
\node[s,draw=omExo,fill=omExo!8] (us) at (9.8,0) {strategies,\\research, risk};
\draw[arr] (ch) -- (ix); \draw[arr] (vx) -- (bb); \draw[arr] (ix) -- (st); \draw[arr] (bb) -- (st); \draw[arr] (st) -- (us);
\end{tikzpicture}

\omcaption{A crypto firm's two data pipelines: blocks from its own nodes through an \omterm{def:m3:crypto-data-and-infrastructure:log}{indexer}, and venues' feeds through a book builder
that checks sequence and checksums, both into one normalised store that keeps the source's timestamp and the time of receipt.
Schematic.}\label{fig:m3:crypto-data-and-infrastructure:pipes}
\end{omfigure}

\section{Venue market data: snapshots, deltas and checksums}

\begin{definition}[Snapshot-and-delta feed, book checksum]\label{def:m3:crypto-data-and-infrastructure:feed}
A \emph{snapshot-and-delta feed}\index{snapshot-and-delta feed} publishes a venue's market-by-price book (One Quant Book 1, chapter 28)
as a snapshot of all levels, with an update identifier, followed by a stream of changes to individual levels numbered with the same
identifiers. A \emph{book checksum}\index{book checksum} is a value the venue computes over the top levels of its book after each change
and publishes with it, so that a client can verify that its copy is identical.
\end{definition}

\begin{dated}{2026-09}{Rebuilding a book: one venue's procedure}\label{dat:m3:crypto-data-and-infrastructure:binance}
Binance's documentation for a local order book instructs the client to buffer the depth stream, fetch a snapshot, discard buffered events
whose last update id $u$ is at most the snapshot's \texttt{lastUpdateId}, and then, for each event: ignore it if $u$ is below the book's id;
if its first update id $U$ exceeds the book's id plus one, discard the book and restart, since events were missed; otherwise set each level to
its new quantity, removing it if the quantity is zero, and set the book's id to $u$. A single websocket connection is valid for 24 hours, and the
server pings every 20 seconds.
\end{dated}

\begin{dated}{2026-09}{A book checksum}\label{dat:m3:crypto-data-and-infrastructure:kraken}
Kraken's documentation for its version-2 websocket book computes the checksum over the top 10 price levels whatever the subscription depth: for asks
from low to high and then bids from high to low, each price and quantity is written without its decimal point and leading zeros, the strings are
concatenated, and the CRC32 of the result, as an unsigned 32-bit integer, is compared with the message's checksum field. Verification is optional.
OKX went the other way: on 23 June 2026 it deprecated the checksum of its order-book channels, still sent but fixed at zero, and told clients to check
continuity with the messages' sequence identifiers instead.
\end{dated}

\begin{omfigure}
\centering
\begin{tikzpicture}[font=\small,
  s/.style={draw,thick,rounded corners=2pt,align=center,minimum height=0.9cm,text width=2.4cm},
  arr/.style={-{Stealth[length=2.2mm]},thick}]
\node[s,draw=omDef,fill=omDef!8] (a) at (0,0) {1.\ subscribe,\\buffer deltas};
\node[s,draw=omThm,fill=omThm!8] (b) at (3.1,0) {2.\ fetch snapshot\\(id $n$)};
\node[s,draw=omProp,fill=omProp!8] (c) at (6.2,0) {3.\ drop buffered\\deltas with $u \le n$};
\node[s,draw=omExo,fill=omExo!8] (d) at (9.3,0) {4.\ apply in order;\\id becomes $u$};
\node[s,draw=gray,fill=gray!8] (e) at (9.3,-2.0) {gap: next $U > \text{id} + 1$};
\draw[arr] (a) -- (b); \draw[arr] (b) -- (c); \draw[arr] (c) -- (d); \draw[arr] (d) -- (e);
\draw[arr,dashed] (e.west) -| node[below,pos=0.25,font=\footnotesize,align=center] {discard the book, flag stale,\\resynchronise} (b.south);
\end{tikzpicture}

\omcaption{Rebuilding a book from a \omterm{def:m3:crypto-data-and-infrastructure:feed}{snapshot-and-delta feed}, following the procedure of \cref{dat:m3:crypto-data-and-infrastructure:binance}: buffer, snapshot,
discard what the snapshot already contains, apply in order, and go back to the snapshot on any gap. The book is flagged stale from the gap until the new snapshot
is applied. Schematic.}\label{fig:m3:crypto-data-and-infrastructure:sync}
\end{omfigure}

Update identifiers catch losses on a venue that publishes them; checksums catch what identifiers cannot, such as a client's own bug in applying a
change, and are the only protection on a feed without identifiers.

\begin{proposition}[What a lost message costs]\label{prop:m3:crypto-data-and-infrastructure:loss}
Let a book of $K$ active levels change by absolute-quantity updates, each touching one level chosen uniformly, and let each message be lost with
probability $p$, independently. A client that applies what it receives without checking is wrong after a loss until the lost level is next updated:
a geometric number of messages with mean about $K$, so that it is silently wrong a fraction about $pK$ of the time for small $p$. A client that detects
the gap and needs $R$ messages to resynchronise is never silently wrong and is knowingly stale a fraction about $pR$ of the time.
\end{proposition}

\begin{proof}
Each later message updates the lost level with probability $1/K$, so the wait is geometric with mean $K$. Losses arrive at rate $p$ per message and, for
small $p$, rarely overlap, so the fractions are the rate times the durations.
\end{proof}

\begin{omfigure}
\begin{tikzpicture}
\begin{axis}[width=11cm,height=5cm,xmode=log,ymode=log,
  xlabel={message loss rate, \%},ylabel={share of time, \%},
  tick label style={font=\small},label style={font=\small},
  xmin=0.008,xmax=1.25,ymin=0.08,ymax=50,grid=major,grid style={gray!25},
  xtick={0.01,0.1,1},xticklabels={0.01,0.1,1},ytick={0.1,1,10},yticklabels={0.1,1,10},
  legend columns=2,legend style={font=\footnotesize,at={(0.5,-0.3)},anchor=north,draw=none,fill=none,/tikz/every even column/.append style={column sep=8pt}},
  table/col sep=comma]
\addplot[omThm,very thick,mark=*,mark size=1.5pt] table[x=p,y=silent]{figdata/markets-3/26-crypto-data-and-infrastructure/loss.csv};\addlegendentry{silently wrong (no checks)}
\addplot[omDef,very thick,dashed,mark=square*,mark size=1.5pt] table[x=p,y=stale]{figdata/markets-3/26-crypto-data-and-infrastructure/loss.csv};\addlegendentry{known stale (gap detection, resync)}
\end{axis}
\end{tikzpicture}

\omcaption{A locally built book of 20 active levels under message loss (\cref{prop:m3:crypto-data-and-infrastructure:loss}). Without checks it is wrong,
without knowing it, about $20p$ of the time; with gap detection it is never silently wrong but spends about $50p$ of the time resynchronising when a snapshot
takes 50 messages to arrive. Simulated feed of 100\,000 messages. Data: the chapter's tutorial.}\label{fig:m3:crypto-data-and-infrastructure:loss}
\end{omfigure}

The careful client is stale more of the time than the careless one is wrong, and that is the point: a stale book that knows it is stale can pull its quotes
and mark its data as missing; a wrong book trades and records falsehoods. Purchased histories must be checked the same way: a vendor's dataset of
incremental book updates is only as good as its recorder's handling of gaps, and a crossed or locked book in the data is the first sign.

\section{Timestamps and what they mean}

Every crypto data record should carry two times: the venue's (exchange timestamp) and the firm's (receive timestamp), as in One Quant Book 1, chapter 28.

\begin{dated}{2026-09}{Timestamps in a vendor's data}\label{dat:m3:crypto-data-and-infrastructure:tardis}
Tardis.dev's downloadable datasets give each record a \texttt{timestamp}, provided by the exchange in microseconds since the epoch (with the local time as a
fallback when the exchange provides none), and a \texttt{local\_timestamp}, the message's arrival time; its incremental book data are collected from the
exchanges' real-time websocket feeds, and files may begin with updates that precede the first snapshot, which must be skipped.
\end{dated}

The exchange's timestamp says when the venue says it happened, at the venue's clock resolution and sometimes when the message was sent rather than when the
book changed. The receive timestamp says when the firm learned it, including the network path. Their difference is the feed's latency, which varies with the
firm's location and the venue's load and jumps exactly when the market is busiest. Research must use the receive timestamp to decide what a strategy could have
known, and the exchange timestamp to align venues; a backtest that uses exchange timestamps for decisions assumes the firm sat inside the matching engine. On
chains there is a third time, the block's, and a fourth, when the firm's node saw the transaction in the \omterm{def:m3:blockchains-for-traders:mempool}{mempool}.

\section{Where the matching engines sit}

\begin{definition}[Cloud region]\label{def:m3:crypto-data-and-infrastructure:region}
A \emph{cloud region}\index{cloud region} is a geographic area in which a cloud provider operates a cluster of data centres, within which network latency between
machines is low.
\end{definition}

Unlike the exchanges of One Quant Book 1, some crypto venues run their systems in public clouds rather than in data centres of their own (one venue's
documentation lists an API endpoint named for a cloud provider, \cref{dat:m3:getting-access-crypto-venues:endpoints}), and a firm that wants to be close to such a
venue rents machines in the same region and, where the provider allows, the same availability zone. Latency to a venue becomes a question of which
region it runs in, which is sometimes published and sometimes inferred by measuring round-trip times from each region; and a firm trading several venues in
different regions faces the same geography as the multi-venue firms of the equity and futures markets, with distances between regions instead of between
buildings.

\begin{dated}{2026-09}{A venue's region}\label{dat:m3:crypto-data-and-infrastructure:region}
Glassnode's research, reported by CoinDesk on 30 March 2026, found Hyperliquid's 24 \omterm{def:m3:blockchains-for-traders:pow}{validators} clustered in Amazon Web Services' Tokyo region, across
several availability zones, so that orders from Tokyo reached them in 2 to 3 milliseconds and orders from Europe in more than 200.
\end{dated}

\section{Tutorial: rebuilding a book and losing a message}\label{tut:m3:crypto-data-and-infrastructure:book}

\begin{tutorial}
\textbf{Goal.} Rebuild an order book from a snapshot and a stream of update-id-numbered deltas, detect gaps and resynchronise, verify the book against a CRC32
checksum, and measure what message loss does. \textbf{End state:} \cref{fig:m3:crypto-data-and-infrastructure:loss} and the numbers of the weekend problem.
\begin{tutsteps}
\item \textbf{The update rule.} Ignore old events, detect gaps, apply absolute quantities.
\omcode{code/firm/wsbook/cpp/firm_wsbook.hpp}{37}{48}{The C++20 update rule: ignore, detect a gap, or apply.}\label{lst:m3:crypto-data-and-infrastructure:apply}
\item \textbf{The checksum.} Ten best asks, then ten best bids, as integer strings; CRC32.
\omcode{code/firm/wsbook/rust/src/lib.rs}{65}{74}{The Rust checksum over the top of the book.}\label{lst:m3:crypto-data-and-infrastructure:checksum}
\item \textbf{Run} \texttt{simulate(p)} for loss rates from 0.01\% to 1\% and \texttt{fig\_data.py}.
\end{tutsteps}
\textbf{What to change next.} Add a checksum policy (detect a divergence at the next message, then resync) and compare it with gap detection; let losses come in
bursts rather than independently; replay a vendor's file and count crossed books.
\end{tutorial}

\section{Build: the websocket book builder}\label{bld:m3:crypto-data-and-infrastructure:wsbook}

\begin{build}
\bfield{Purpose} Every price the miniature firm uses on a crypto venue comes from a book it built itself from the venue's feed; the builder must never present a
wrong book as right, and must say when it is stale. It feeds the arbitrage scanner of \cref{ch:m3:spot-markets}.
\bfield{Interface} \texttt{Book} with \texttt{snapshot(update\_id, bids, asks)}, \texttt{apply(first, last, bids, asks)} returning ignored, applied or gap,
\texttt{top(n)} and \texttt{checksum()}; \texttt{crc32}. C++20 header \texttt{firm\_wsbook.hpp}, Rust crate \texttt{firm\_wsbook}, Python reference.
\bfield{Rules} Integer ticks and lots; a gap leaves the book unsynced until a snapshot; the checksum computed exactly as the venue does; every update keeps both
timestamps.
\bfield{Acceptance tests} \texttt{code/firm/wsbook/tests/}, \texttt{cpp/firm\_wsbook\_test.cpp} and the Rust crate's tests: ignoring, applying and gap
detection; staying unsynced; the checksum against zlib's CRC32 and the standard check value.
\bfield{Stretch} Venue adapters (decimal strings, depth limits of snapshots); buffered resynchronisation without dropping events; burst-loss testing; a
lock-free publication of the book to strategies.
\end{build}

\begin{omsources}
\item Ethereum.org developer documentation, ``Nodes and clients'' (repository ethereum/ethereum-org-website). Accessed September 2026.
\item Binance, spot WebSocket streams documentation (repository binance/binance-spot-api-docs); Kraken, WebSocket v2 book checksum guide. Accessed September 2026.
\item Tardis.dev, downloadable CSV files: data types. Accessed September 2026.
\item OKX, API v5 change log, entry of 23 June 2026; CoinDesk, ``Hyperliquid traders in Tokyo get 200-millisecond edge, Glassnode research shows'', 30 March 2026.
\end{omsources}

\section{Exercises}

\begin{exercise}[$\star$]\label{exo:m3:crypto-data-and-infrastructure:1}
A book's id is 100. What does the client do with events $(U, u) = (95, 99)$, $(101, 103)$ and then $(105, 106)$?
\end{exercise}

\begin{exercise}[$\star$]\label{exo:m3:crypto-data-and-infrastructure:2}
When does a firm need an \omterm{def:m3:crypto-data-and-infrastructure:node}{archive node} rather than a \omterm{def:m3:crypto-data-and-infrastructure:node}{full node}?
\end{exercise}

\begin{exercise}[$\star$]\label{exo:m3:crypto-data-and-infrastructure:3}
Write the checksum string for asks (1001, 4), (1002, 6) and bids (999, 5), (998, 7).
\end{exercise}

\begin{exercise}[$\star\star$]\label{exo:m3:crypto-data-and-infrastructure:4}
A feed loses one message in 1\,000 and the book has 20 active levels. For what share of the time is an unchecked book wrong? How long is a typical error?
\end{exercise}

\begin{exercise}[$\star\star$]\label{exo:m3:crypto-data-and-infrastructure:5}
Why should research use receive timestamps to decide what a strategy knew?
\end{exercise}

\begin{exercise}[$\star\star$]\label{exo:m3:crypto-data-and-infrastructure:6}
What must an \omterm{def:m3:crypto-data-and-infrastructure:log}{indexer} do when the chain reorganises two blocks?
\end{exercise}

\begin{exercise}[$\star\star\star$]\label{exo:m3:crypto-data-and-infrastructure:7}
\emph{Coding.} With \texttt{simulate}, find the resynchronisation time at which the careful client is stale as often as the careless one is wrong, at a loss rate of
0.1\%.
\end{exercise}

\begin{exercise}[$\star\star\star$]\label{exo:m3:crypto-data-and-infrastructure:8}
\emph{Find the flaw.} ``Our data vendor records every exchange's feed, so its order books are correct.''
\end{exercise}

\section{Problem: The Missing Delta}

\begin{problem}[{Weekend problem --- how wrong is a book that lost a message?}]\label{pb:m3:crypto-data-and-infrastructure:1}
A firm builds books from a venue's feed of 20 active levels near the top. It loses one message in 1\,000. A resynchronisation from a snapshot takes about 50
messages.

\medskip\noindent\textbf{Part I --- The careless client.}
\begin{enumerate}
\item For what share of messages is its book wrong?
\item How long does a typical error last, in messages?
\item Why is that close to the number of levels?
\item What can a wrong book do to a quoting strategy?
\item Why might a lost deletion be worse than a lost change of quantity?
\end{enumerate}

\medskip\noindent\textbf{Part II --- The careful client.}
\begin{enumerate}[resume]
\item For what share of messages is its book flagged stale?
\item Why is that larger than the careless client's error share?
\item What does the strategy do while the book is stale?
\item How would buffering events during resynchronisation shorten the stale time?
\item What would checksums add on a venue that publishes update identifiers?
\end{enumerate}

\medskip\noindent\textbf{Part III --- Data for research.}
\begin{enumerate}[resume]
\item How would you detect lost messages in a purchased dataset?
\item What does a crossed book in the data indicate?
\item Which timestamp should a backtest use for decisions?
\item How do you align two venues' books in time?
\item What should the dataset record about resynchronisations?
\end{enumerate}

\medskip\noindent\textbf{Part IV --- Judgement.}
\begin{enumerate}[resume]
\item Is a stale book better than a wrong one?
\item When is a provider's node enough, and when must a firm run its own?
\item What makes crypto data harder than equity data, and what makes it easier?
\item State the \emph{named result}: the share of time the book is wrong without checks and stale with resynchronisation, and the expected duration of an error.
\item In one sentence: what is a local order book?
\end{enumerate}
\end{problem}

\section{Interview questions}

\begin{interviewq}[$\star$ \iqroles{developer}]\label{iq:m3:crypto-data-and-infrastructure:1}
Describe how to maintain a local order book from a snapshot and a delta stream.
\end{interviewq}

\begin{interviewq}[$\star$ \iqroles{researcher}]\label{iq:m3:crypto-data-and-infrastructure:2}
What is the difference between an exchange timestamp and a receive timestamp, and when do you use each?
\end{interviewq}

\begin{interviewq}[$\star\star$ \iqroles{developer}]\label{iq:m3:crypto-data-and-infrastructure:3}
How would you design an \omterm{def:m3:crypto-data-and-infrastructure:log}{indexer} for every swap on a \omterm{def:m3:automated-market-makers:amm}{decentralised exchange}, robust to reorganisations?
\end{interviewq}

\begin{interviewq}[$\star\star$ \iqroles{developer}]\label{iq:m3:crypto-data-and-infrastructure:4}
Your \omterm{def:m3:crypto-data-and-infrastructure:feed}{book checksum} fails once an hour on one venue. How do you find out why?
\end{interviewq}

\begin{interviewq}[$\star\star$ \iqroles{researcher}]\label{iq:m3:crypto-data-and-infrastructure:5}
How would you measure the latency of your feed from a venue running in a public cloud?
\end{interviewq}

\begin{interviewq}[$\star\star\star$ \iqroles{developer}]\label{iq:m3:crypto-data-and-infrastructure:6}
Design the market-data layer for a firm trading on ten crypto venues and three chains.
\end{interviewq}
